% !TeX program = xelatex
% Perihelion advance shown by two successive osculating ellipses.
% Compile: xelatex perihelion_precession_tikz.tex
\documentclass[tikz,border=6pt]{standalone}
\usepackage{amsmath}
\usepackage[UTF8,fontset=fandol]{ctex}
\usepackage{fix-cm}
\usetikzlibrary{arrows.meta}

\begin{document}
\begin{tikzpicture}[>=Stealth,thick,every node/.style={font=\fontsize{16}{19}\selectfont}]
  % Ellipse a=4, b=3.2, focus at the origin: centre (-2.4,0).
  \draw[blue!70!black,dashed,line width=1.15pt] (-2.4,0) ellipse (4 and 3.2);
  \begin{scope}[rotate=18]
    \draw[red,line width=1.25pt] (-2.4,0) ellipse (4 and 3.2);
  \end{scope}

  \shade[inner color=yellow!20,outer color=orange!85] (0,0) circle (.32);
  \node[below left] at (-.12,-.28) {$M_\odot$};

  \coordinate (Pzero) at (1.6,0);
  \coordinate (Pone) at ({1.6*cos(18)},{1.6*sin(18)});
  \fill[blue!70!black] (Pzero) circle (2.2pt);
  \fill[red] (Pone) circle (2.2pt);
  \draw[blue!70!black,thin] (0,0) -- (Pzero);
  \draw[red,thin] (0,0) -- (Pone);
  \draw[green!45!black,->,line width=1.2pt] (.95,0) arc[start angle=0,end angle=18,radius=.95];
  \node[green!40!black] at (2.4,0.2) {$\Delta\phi$};

  \node[blue!70!black,anchor=west] at (2.35,-0.6) {第 $n$ 圈};
  \draw[blue!70!black,->] (2.45,-0.45) -- (1.75,-0.09);
  \node[red,anchor=west] at (2.35,1.35) {第 $n+1$ 圈};
  \draw[red,->] (2.4,1.2) -- (1.6,0.6);
  \node[align=left,anchor=west] at (3.0,-2.5)
    {轨道近似为开普勒椭圆，\\但近日点缓慢进动};
\end{tikzpicture}
\end{document}
